\documentclass[11pt]{article}
\usepackage[utf8]{inputenc}
\usepackage{amsmath,amssymb}
\usepackage[margin=1in]{geometry}
\usepackage{hyperref}
\emergencystretch=3em

\title{The exponential tree series in $d$ dimensions}
\author{Claude, with Daniel Murfet}
\date{2026-09-06}

\begin{document}
\maketitle
\sloppy

\begin{abstract}
The proof of \texttt{thm:TaylorTree} begins by writing $\xi=\xi(0)+J$ and expanding
$e^{\beta\sqrt n\,u^kJ(u)}$ in its exponential series. For any continuous $J$ on the box the series
may be integrated term by term:
$\int u^he^{-\beta nu^{2k}+\beta\sqrt n\,u^k(a+J)}\,du=\sum_{p\ge0}\frac{\beta^p}{p!}\int u^h(\sqrt n\,u^kJ)^pf(\sqrt n\,u^k)\,du$
as a convergent series (\texttt{boxIntegralXi\_hasSum}). Zero \texttt{sorry}/\texttt{axiom}.
\end{abstract}

\section{Setting}
On the compact box every factor is bounded: $u^h\le\prod b^{h_i}$, $\sqrt n\,u^k\le\sqrt n\prod b^{k_i}$,
$|J|\le C_J$ (continuous on the closed box), $f\le E$. The $p$-th layer is therefore dominated by the
constant $H_bE\,(\beta K_bC_J)^p/p!$, whose sum is $H_bE\,e^{\beta K_bC_J}$, integrable against the
finite box measure; dominated convergence for series does the rest.

\section{The things formalised}
{\small
\begin{verbatim}
noncomputable def boxIntegralXi (β a b c : ℝ) {d : ℕ} (k h : Fin d → ℕ)
    (J : (Fin d → ℝ) → ℝ) : ℝ :=
  ∫ u, (∏ i, u i ^ h i)
    * Real.exp (-β * (c * ∏ i, u i ^ k i) ^ 2
        + β * (c * ∏ i, u i ^ k i) * (a + J u)) ∂(boxMeasure b d)
noncomputable def treeLayer (β a b c : ℝ) {d : ℕ} (k h : Fin d → ℕ)
    (J : (Fin d → ℝ) → ℝ) (p : ℕ) : ℝ :=
  β ^ p / p.factorial * ∫ u, (∏ i, u i ^ h i) * ((c * ∏ i, u i ^ k i) * J u) ^ p
    * quadKernel β a (c * ∏ i, u i ^ k i) ∂(boxMeasure b d)
theorem boxIntegralXi_hasSum (β a b c : ℝ) (hβ : 0 < β) (hb : 0 < b) (hc : 0 ≤ c)
    {d : ℕ} (k h : Fin d → ℕ) (J : (Fin d → ℝ) → ℝ) (hJ : Continuous J) :
    HasSum (treeLayer β a b c k h J) (boxIntegralXi β a b c k h J)
\end{verbatim}
}

\section{Proof strategy}
Unit~26's one-dimensional idiom (\texttt{hasSum\_integral\_of\_dominated\_convergence}) transferred
verbatim: measurability from continuity; the almost-everywhere bound from
\texttt{ae\_restrict\_mem} on the box (the box measure is a restriction, unit~54) and the compact bound
\texttt{IsCompact.exists\_bound\_of\_continuousOn}; summability and integrability of the constant
bound; and the pointwise exponential series \texttt{NormedSpace.expSeries\_div\_hasSum\_exp} with
$e^{-\beta s^2+\beta s(a+J)}=f(s)e^{\beta sJ}$.

\section{Roadblocks and resolutions}
\begin{itemize}
\item \texttt{simp} rewrote membership of the constant point in the closed box into a pointwise
inequality of functions; \texttt{Set.mem\_univ\_pi} and an explicit witness instead.
\end{itemize}

\section{What was Mathlib, what was new}
Mathlib supplied the series interchange (\texttt{hasSum\_integral\_of\_dominated\_convergence}), the
compact bound (\texttt{IsCompact.exists\_bound\_of\_continuousOn} with \texttt{isCompact\_univ\_pi}),
the summable majorant (\texttt{Real.summable\_pow\_div\_factorial}), the exponential series
(\texttt{NormedSpace.expSeries\_div\_hasSum\_exp}) and \texttt{Finset.prod\_le\_prod}. New: the
$d$-dimensional tree series.

\section{Lessons}
The compactness of the chart box makes the analytic step of the Taylor tree trivial for any continuous
$J$; the difficulty of the theorem lies entirely in the combinatorics of expanding $J^p$ and in
summing the resulting tree, not in the interchange of sum and integral.

\section{Outcome}
The Taylor tree's outer expansion is formalised in $d$ dimensions. Each layer is, for polynomial $J$,
a finite combination of the tree terms of unit~60, all of whose leading exponents lie in $\Lambda(h,k)$.
\end{document}
